%%% The main file. It contains definitions of basic parameters and includes all other parts.

%% Settings for single-side (simplex) printing
% Margins: left 40mm, right 25mm, top and bottom 25mm
% (but beware, LaTeX adds 1in implicitly)
\documentclass[12pt,a4paper]{report}
\setlength\textwidth{145mm}
\setlength\textheight{247mm}
\setlength\oddsidemargin{15mm}
\setlength\evensidemargin{15mm}
\setlength\topmargin{0mm}
\setlength\headsep{0mm}
\setlength\headheight{0mm}
% \openright makes the following text appear on a right-hand page
\let\openright=\clearpage

%% Settings for two-sided (duplex) printing
% \documentclass[12pt,a4paper,twoside,openright]{report}
% \setlength\textwidth{145mm}
% \setlength\textheight{247mm}
% \setlength\oddsidemargin{14.2mm}
% \setlength\evensidemargin{0mm}
% \setlength\topmargin{0mm}
% \setlength\headsep{0mm}
% \setlength\headheight{0mm}
% \let\openright=\cleardoublepage

%% Character encoding: usually latin2, cp1250 or utf8:
\usepackage[utf8]{inputenc}

%% Further useful packages (included in most LaTeX distributions)
\usepackage{amsmath}        % extensions for typesetting of math
\usepackage{amsfonts}       % math fonts
\usepackage{amsthm}         % theorems, definitions, etc.
\usepackage{bbding}         % various symbols (squares, asterisks, scissors, ...)
\usepackage{bm}             % boldface symbols (\bm)
\usepackage{graphicx}       % embedding of pictures
\usepackage{fancyvrb}       % improved verbatim environment
\usepackage{natbib}         % citation style AUTHOR (YEAR), or AUTHOR [NUMBER]
\usepackage[nottoc]{tocbibind} % makes sure that bibliography and the lists
			    % of figures/tables are included in the table
			    % of contents
\usepackage{dcolumn}        % improved alignment of table columns
\usepackage{booktabs}       % improved horizontal lines in tables
\usepackage{paralist}       % improved enumerate and itemize
\usepackage[usenames]{xcolor}  % typesetting in color

%%% Basic information on the thesis

% Thesis title in English (exactly as in the formal assignment)
\def\ThesisTitle{Multicast routing}

% Author of the thesis
\def\ThesisAuthor{Ondřej Hlavatý}

% Year when the thesis is submitted
\def\YearSubmitted{2016}

% Name of the department or institute, where the work was officially assigned
% (according to the Organizational Structure of MFF UK in English,
% or a full name of a department outside MFF)
\def\Department{Department of Applied Mathematics}

% Is it a department (katedra), or an institute (ústav)?
\def\DeptType{Department}

% Thesis supervisor: name, surname and titles
\def\Supervisor{Mgr. Martin Mareš, Ph.D.}

% Supervisor's department (again according to Organizational structure of MFF)
\let\SupervisorsDepartment\Department

% Study programme and specialization
\def\StudyProgramme{Computer Science}
\def\StudyBranch{IOI}

% An optional dedication: you can thank whomever you wish (your supervisor,
% consultant, a person who lent the software, etc.)
\def\Dedication{%
	\TODO Dedication.
}

% Abstract (recommended length around 80-200 words; this is not a copy of your thesis assignment!)
\def\Abstract{%
There are multiple approaches to multicast routing in IPv4 networks. BIDIR-PIM
is the only protocol in PIM family that doesn't have to hold source specific
information. Multicast routing kernel interface is complicated and archaic. }

% 3 to 5 keywords (recommended), each enclosed in curly braces
\def\Keywords{%
	{multicast} {routing} {bird} {unix}
}

%% The hyperref package for clickable links in PDF and also for storing
%% metadata to PDF (including the table of contents).
\usepackage[pdftex,unicode]{hyperref}   % Must follow all other packages
\hypersetup{breaklinks=true}
\hypersetup{pdftitle={\ThesisTitle}}
\hypersetup{pdfauthor={\ThesisAuthor}}
\hypersetup{pdfkeywords=\Keywords}
\hypersetup{urlcolor=blue}

% Definitions of macros (see description inside)
%%% This file contains definitions of various useful macros and environments %%%
%%% Please add more macros here instead of cluttering other files with them. %%%

%%% Minor tweaks of style

% These macros employ a little dirty trick to convince LaTeX to typeset
% chapter headings sanely, without lots of empty space above them.
% Feel free to ignore.
\makeatletter
\def\@makechapterhead#1{
  {\parindent \z@ \raggedright \normalfont
   \Huge\bfseries \thechapter. #1
   \par\nobreak
   \vskip 20\p@
}}
\def\@makeschapterhead#1{
  {\parindent \z@ \raggedright \normalfont
   \Huge\bfseries #1
   \par\nobreak
   \vskip 20\p@
}}
\makeatother

% This macro defines a chapter, which is not numbered, but is included
% in the table of contents.
\def\chapwithtoc#1{
\chapter*{#1}
\addcontentsline{toc}{chapter}{#1}
}

% Draw black "slugs" whenever a line overflows, so that we can spot it easily.
\overfullrule=1mm

%%% Macros for definitions, theorems, claims, examples, ... (requires amsthm package)

\theoremstyle{plain}
\newtheorem{thm}{Theorem}
\newtheorem{lemma}[thm]{Lemma}
\newtheorem{claim}[thm]{Claim}

\theoremstyle{plain}
\newtheorem{defn}{Definition}

\theoremstyle{remark}
\newtheorem*{cor}{Corollary}
\newtheorem*{rem}{Remark}
\newtheorem*{example}{Example}

%%% An environment for proofs

%%% FIXME %%% \newenvironment{proof}{
%%% FIXME %%%   \par\medskip\noindent
%%% FIXME %%%   \textit{Proof}.
%%% FIXME %%% }{
%%% FIXME %%% \newline
%%% FIXME %%% \rightline{$\square$}  % or \SquareCastShadowBottomRight from bbding package
%%% FIXME %%% }

%%% An environment for typesetting of program code and input/output
%%% of programs. (Requires the fancyvrb package -- fancy verbatim.)

\DefineVerbatimEnvironment{code}{Verbatim}{fontsize=\small, frame=single}

%%% The field of all real and natural numbers
\newcommand{\R}{\mathbb{R}}
\newcommand{\N}{\mathbb{N}}

%%% Useful operators for statistics and probability
\DeclareMathOperator{\pr}{\textsf{P}}
\DeclareMathOperator{\E}{\textsf{E}\,}
\DeclareMathOperator{\var}{\textrm{var}}
\DeclareMathOperator{\sd}{\textrm{sd}}

%%% Transposition of a vector/matrix
\newcommand{\T}[1]{#1^\top}

%%% Various math goodies
\newcommand{\goto}{\rightarrow}
\newcommand{\gotop}{\stackrel{P}{\longrightarrow}}
\newcommand{\maon}[1]{o(n^{#1})}
\newcommand{\abs}[1]{\left|{#1}\right|}
\newcommand{\dint}{\int_0^\tau\!\!\int_0^\tau}
\newcommand{\isqr}[1]{\frac{1}{\sqrt{#1}}}

%%% Various table goodies
\newcommand{\pulrad}[1]{\raisebox{1.5ex}[0pt]{#1}}
\newcommand{\mc}[1]{\multicolumn{1}{c}{#1}}

\let\ttt\texttt
\newcommand{\TODO}[1]{\textbf{TODO #1}}
\newcommand{\prefix}[1]{\ttt{#1}}
\newcommand{\ip}[1]{\ttt{#1}}
\newcommand{\mac}[1]{\ttt{#1}}
\newcommand{\rfc}[1]{RFC~#1 \cite{rfc#1}}
\newcommand{\NULL}{\ttt{NULL}}



% Title page and various mandatory informational pages
\begin{document}
%%% Title page of the thesis and other mandatory pages

%%% Title page of the thesis

\pagestyle{empty}
\hypersetup{pageanchor=false}
\begin{center}

\large

Charles University in Prague

\medskip

Faculty of Mathematics and Physics

\vfill

{\bf\Large BACHELOR THESIS}

\vfill

\centerline{\mbox{\includegraphics[width=60mm]{img/logo.pdf}}}

\vfill
\vspace{5mm}

{\LARGE\ThesisAuthor}

\vspace{15mm}

{\LARGE\bfseries\ThesisTitle}

\vfill

\Department

\vfill

\begin{tabular}{rl}

Supervisor of the bachelor thesis: & \Supervisor \\
\noalign{\vspace{2mm}}
Study programme: & \StudyProgramme \\
\noalign{\vspace{2mm}}
Study branch: & \StudyBranch \\
\end{tabular}

\vfill

% Zde doplňte rok
Prague \YearSubmitted

\end{center}

\newpage

%%% Here should be a bound sheet included -- a signed copy of the "bachelor
%%% thesis assignment". This assignment is NOT a part of the electronic
%%% version of the thesis. DO NOT SCAN.

%%% A page with a solemn declaration to the bachelor thesis

\openright
\hypersetup{pageanchor=true}
\pagestyle{plain}
\pagenumbering{roman}
\vglue 0pt plus 1fill

\noindent
I declare that I carried out this bachelor thesis independently, and only with the cited
sources, literature and other professional sources.

\medskip\noindent
I understand that my work relates to the rights and obligations under the Act No.~121/2000 Sb.,
the Copyright Act, as amended, in particular the fact that the Charles
University in Prague has the right to conclude a license agreement on the use of this
work as a school work pursuant to Section 60 subsection 1 of the Copyright Act.

\vspace{10mm}

\hbox{\hbox to 0.5\hsize{%
In ........ date ............	% FIXME!
\hss}\hbox to 0.5\hsize{%
signature of the author
\hss}}

\vspace{20mm}
\newpage

%%% Mandatory information page of the thesis

\openright

\vbox to 0.5\vsize{
\setlength\parindent{0mm}
\setlength\parskip{5mm}

Title:
\ThesisTitle

Author:
\ThesisAuthor

\DeptType:
\Department

Supervisor:
\Supervisor, \SupervisorsDepartment

Abstract:
\Abstract

Keywords:
\Keywords

\vss}

\newpage

%%% Dedication

\openright

\noindent
\Dedication

\newpage

\openright
\pagestyle{plain}
\pagenumbering{arabic}
\setcounter{page}{1}


%%% A page with automatically generated table of contents of the bachelor thesis

\tableofcontents

%%% Each chapter is kept in a separate file
\chapter*{Introduction}
\addcontentsline{toc}{chapter}{Introduction}

There are multiple ways how to use a computer network to communicate between two
or more machines. In the most common technology, based on the Internet Protocol
(IP), one can send a message (called a packet) from one machine, directed to
another one. That is usually called a \emph{unicast} and is by far the most well-known
method. Or one can send a packet to all machines in some network -- that is
called a \emph{broadcast}. However, there are applications where both of these
are suboptimal. A good example can be conference telephony, games or video
streaming.
\nomenclature{IP}{Internet Protocol}

When a public television company broadcasts the final match of Olympic games, it can
have millions of receivers. The company does not want to send every packet in
millions of copies by unicast. However, it does not want to flood every
computer connected to the Internet by broadcast either. That is why \emph{multicast} has
been invented. Multicast is a general way how to transmit traffic from a group of
senders to a group of receivers.

In order to deliver packets between networks, the network itself has to do some
computations. One of them is routing -- collaboratively finding an optimal path
for a packet, for some definition of optimal.
The world of unicast routing is well inhabited. There are lots of unicast routing
protocols based on completely different approaches. When it comes to routing
multicasts, it is much different.

Multicast itself was added as an extension to IPv4 and it requires special
handling in network hardware. When just one router on the way does not support
it, multicast will not work. As such it suffered from a lack of support for a long
time. And because this support was often broken, it is still not widely used.

Supporting multicast on the side of the receiver is easy, and one can receive
multicast in all current major operating systems. For example in Unix-like operating systems, it is
a~matter of a few syscalls. But routing multicast traffic is a~different thing
-- one needs to run a routing daemon that communicates with hosts and other
routers, speaks at least two different network protocols and controls the
kernel's routing tables. There already are such daemons that are traditionally
used: \texttt{mrouted} and \texttt{pimd}.

The purpose of this thesis is to explain how multicast works, introduce the
reader to the world of multicast routing and implement one of existing
multicast routing protocols. It would be possible to implement a standalone
routing daemon, but its chance to catch on would be small.

There once was a student project at MFF UK, that grew to worldwide usage -- the BIRD
internet routing daemon. It currently supports a handful of unicast routing
protocols, that can run at the same time. There are use-cases in which having closer communication
between unicast and multicast routing can be an advantage. Let the primary goal of
this thesis be to teach BIRD multicast routing.%
\nomenclature{BIRD}{BIRD Internet Routing Daemon}

The protocol chosen for this thesis is PIM-BIDIR. It is a somehow simplified
protocol from the PIM family because it doesn't need to hold state information
based on traffic source. That makes it less demanding on network hardware.

\chapter{Multicast in IPv4}

IP multicast was added as an additional addressing model. The first
references can be found in \rfc{966}. This RFC defines multicast traffic
with many responsibilities for routers (called ``multicast agents'' at that time)
and it was soon obsoleted and served only as an inspiration. The latest
specification can be found in \rfc{1112}.

In local networks, multicasts are similar to broadcasts. Every packet is to be
transmitted to every host on the same link, and each host can decide according to
the destination address whether or not the packet should be processed or altogether
dropped. In other words, you need neither a multicast routing capable router
nor any special network hardware to use multicast in one network. Differences
are notable on inter-network level. While broadcasts should be either forwarded
to all links or filtered completely by routers, multicasts should have a more
controlled behavior.

\section{Multicast groups}

The first term we need to understand is the \emph{multicast group}. It is represented by an IP
address allocated from a dedicated region \prefix{224.0.0.0/4}. This region is further
divided into blocks, which have additional semantics. From routing point of
view, only the Local Network Control Block (\prefix{224.0.0.0/24}) is
interesting. Multicasts in this group must never be forwarded and they stay
in the boundary of a~local network. This block is often used by unicast
routing protocols.

Host membership in a group is fully optional. Neither the number of members in
a group nor the number of groups a host belongs to is limited and it can be even
zero. A host does not have to belong to a group to send packets into it.
Multicast packets are distributed in the same ``best effort'' manner as unicast ones -- they can
be delivered to all, some or no hosts, and in any order.

There are three levels of conformance with IP multicast. Level 0 means no
conformance. Level 1 requires hosts to be able to send packets to multicast
groups. This level still requires little code in the networking stack, because
sending a~packet to a multicast group is very similar to sending it as an unicast
packet with a destination address set to the group address. While it doesn't make much
sense to send unicast packets destined outside the local network prefix directly
to the link, it is the expected behavior with multicast. Level 2 means fully
conforming, able to send and receive multicast traffic.

\section{Internet Group Management Protocol}

As already said, membership in a multicast group is optional. By default, there
is only one group that every host listens to --- the \ttt{ALL-HOSTS} group with address
\ip{224.0.0.1}. There must be a way for hosts to express group membership.
A~protocol exists for this purpose and it is called IGMP.
\nomenclature{IGMP}{Internet Group Management Protocol}

The most recent version of IGMP is 3. It is described in \rfc{3376}. However, all
versions are backwards compatible. Whenever there is a router which only knows
an older version, all the newer ones stop sending any packets and they only receive and observe.

The purpose of IGMP is to exchange information about group membership between
hosts and routers. Note that it is mandatory to implement IGMP in order to participate in
a multicast group, but until a multicast packet crosses the boundary of a single network, it
is delivered as broadcast to all hosts. As such all hosts receive it without
the routers being involved.

Let us start with IGMP version 1 specified in Appendix I of \rfc{1112}. There
are just two types of messages: \emph{membership query} and \emph{membership report}.
Membership query is sent periodically by routers. It announces to hosts that there
is an IGMPv1 capable router. This query is sent to \ttt{ALL-HOSTS} group with TTL~1 and
therefore is limited to the local network. When a host receives a query, it replies with
a membership report in 10 seconds with a random delay. No host should send
reports for the \ttt{ALL-HOSTS} group, membership there is implicit and mandatory.

The report is sent to the group being reported. As noted above, the host does not have
to send reports in order to listen. Filtering multicast traffic is made by other means.
Sending a~report declares there is someone interested in this group on this
particular link. When, during the random delay, the host receives a~report, it
stops his own timer and remains silent. This reduces the amount of control
traffic. When the host decides to leave a group, it does nothing. If it was the
only host in that group, on the next query no one will report and the routers will
stop forwarding.

IGMPv2 \cite{rfc2236} comes with major changes. First, a router can send
a query to a~single specific group. That means it can keep a timer for every group
separately, and asking for reports does not generate unnecessary traffic in
every other group. Second, the maximum delay between a query and the report
(previously fixed to 10 seconds) is now specified in each query and it is
configurable. A network administrator can tune this variable for his needs.
When the value is decreased, hosts will reply faster and routers can faster stop
forwarding. Giving longer delays can improve stability on very slow or
unreliable links.

The most notable change though is a new message type -- \emph{leave group}. When a host
sends a~report, it sets a flag. Other hosts without the flag remain silent. When
the flagged host leaves a group, it sends the leave group message. Other hosts
react to other hosts' leave in a way similar to a query. When a router hears the leave,
it starts a timer. This timer can be stopped by receiving a~report. Expiring this
timer immediately changes this group state to not having any members. This way
less unnecessary traffic is forwarded.

IGMPv3 comes with additional fields and more complex protocol logic. It is
dedicated to source-specific multicast forwarding and includes messages to manage
a set of sources for every link. Routers should then forward only traffic from
these sources onto this link.

For purpose of this thesis, only IGMPv2 was implemented, as PIM-BIDIR is not
a~source-specific routing protocol.

\section{Link layer}

Even though multicast is specified on the network layer, it needs some level of
support from the link layer. Let us explain how multicast traffic integrates with
the most common combination of IP over Ethernet.

The first problem we need to solve is which MAC address to use for multicast
group addresses. When sending an unicast packet, a mechanism called Address
Resolution Protocol (\emph{ARP}, \rfc{826}) is used to determine the MAC
address. It is done simply by broadcasting questions ``Who has this IP
address?'' and collecting answers and caching them.
\nomenclature{MAC}{Media Access Control}

As the multicast is technically a broadcast on the link layer, it would be possible to
use the broadcast address \mac{FF:FF:FF:FF:FF:FF}.
But that would mean that every multicast packet on the network must be received by all
hosts and examined by their operating systems, because there is no way how to filter
it on the hardware level. For this purpose, every Ethernet address with the
least-significant bit of the first byte set is considered to be
a broadcast address and is to be flooded to all nodes on this network. Network
hardware generally does not make a difference between them, so we can associate
specific meaning to different MAC addresses.

For IPv4 multicast, a region of MAC addresses was allocated. This region starts
at \mac{01:00:5E:00:00:00} and ends with \mac{01:00:5E:7F:FF:FF}. You can see
that this address space is smaller (only 23 bits) than the multicast IP region
(28~bits). To construct the MAC address from the multicast group address, one
puts lower 23 bits of the IP address to the lower bits of the MAC address. This
way, network interfaces can implement a rough filter, that can drop frames the
host knows nothing about. Additional filtering is still needed because of those
missing bits.

A multicast router needs to receive packets for all groups because IGMP reports
are sent there. To enable this behavior, interfaces often have an ``all multicast''
mode, which lets all multicast frames pass through. This mode is also used when
the filter table on the interface is too small to contain all the filters needed. Or
we can put the interface into a so-called promiscuous mode, which passes all frames
to the operating system.

\subsection{IGMP snooping}

What was written above is true according to the specifications, but not really
reflecting the real state of the art. Link layer hardware such as switches often does
make a difference and it uses a technique called \emph{IGMP snooping}. This is described
in \rfc{4541}. This RFC is only informational, because a lot of network hardware with
this functionality was produced before this RFC and it does not behave in any
standardized way. Network switches performing IGMP snooping understand IGMP. They
listen to routers communicating with hosts and they filter multicast traffic on
the link layer.

At the first glance, this seems reasonable, because it stops forwarding packets
to network segments without any receivers, thus saving bandwidth. On the other
hand, this high-level functionality is expected from the network layer, and
reimplementing it in the link layer brings many corner cases, which can prevent
local multicast from working at all. Especially when connecting two multicast
routers with a bridge, you must disable IGMP snooping, because it will filter
out all traffic for which the router is not a local host -- most probably
everything you want to route.

For instance, Linux virtual bridges have IGMP snooping turned on by default.
You can check its status and configure it through special files located in
\ttt{/sys/devices/virtual/{\it netdev\/}/bridge/}.

\def\setsockopt{\ttt{set\-sock\-opt} }
\section{Kernel implementation}

The system interface for multicast is pretty complicated. At least, differences
between Unixes are not big. Let us start with code to receive multicast traffic
in an ordinary application. It is not much harder than opening a socket to
receive unicasts:

\lstinputlisting[firstline=8, lastline=18]{examples/socket.c}

\noindent You might have expected that we would \emph{bind} the socket to the
multicast group address. That is not true. A socket can join an arbitrary
number of groups. Instead, we're binding it to a wildcard address and joining
the group issuing a~\setsockopt call.

One thing to note here. When there is a socket opened which joined a group, the host itself joins the group by means of IGMP. And when the host is
joined, every other multicast-enabled socket will get all packets addressed to that
group. So, you still need to filter those at the transport layer (using ports
for example) or by yourself in the application.

When the router daemon wants to receive all the IGMP packets for all groups, it
would have to join every group possible. That is certainly not the way to go.
IGMPv2 suggests using the IP option \emph{Router Alert}. Every IP packet with this option
set should be passed to all sockets that order it by this \setsockopt call:

\lstinputlisting[firstline=8, lastline=11]{examples/igmp_socket.c}

\noindent There are multiple problems with this. Backward compatibility with
IGMPv1 is broken, and standard does not force hosts to set the option. Some (e.g.,
old Cisco) routers do not do it.

Surprisingly, there is only one way how to get all IGMP packets coming on
some interface. That is opening the IGMP control socket.

\subsection{Multicast kernel interface}
\label{mcast-kernel}

When we want to control kernel multicast routing tables, we need to open
a~control socket. That is a raw IP socket with a type \ttt{IPPROTO\_IGMP}.
Then, a~\setsockopt call is needed:

\begin{lstlisting}[language=c]
int fd = socket(AF_INET, SOCK_RAW, IPPROTO_IGMP);

int v = 1;
setsockopt(fd, IPPROTO_IP, MRT_INIT, &v, sizeof(v));
\end{lstlisting}

\noindent Syscalls \ttt{setsockopt} related to the multicast have option names starting
with \ttt{MRT\_}. These constants are defined in \ttt{<linux/mroute.h>}.
Because they are essentially function calls, allow us to say ``call
\ttt{MRT\_INIT}'' with the meaning of calling \setsockopt with the option
name \ttt{MRT\_INIT}.

There can be only one socket on which the call \ttt{MRT\_INIT} succeeds. That socket
is stored in the kernel and it is used for multiple purposes. By issuing further
\setsockopt calls, one can control multicast routing tables. What is not well
documented is that this is the only socket on a host which can receive all IGMP
packets. Because of this limitation, it is not possible to run more than one
multicast routing daemon.

Let us have a look on how to control the multicast routing. As the first thing
we need to add \emph{virtual interfaces} (VIFs). A VIF can be built on top of
a~real network device, an IPIP tunnel, or it can create a special purpose device.
Multicast routing table entries operate only on subsets of VIFs.
\nomenclature{VIF}{Virtual Interface}

There can be only a very limited number of VIFS. This number is currently
controlled by a kernel compile-time constant \ttt{MAXVIFS}, which defaults to 32. It is not easy
to change it, because it breaks the kernel ABI. Every VIF has an index from the
range $0\dots31$. To add a VIF, we must fill \ttt{struct vifctl} and call
\ttt{MRT\_ADD\_VIF}:

\begin{lstlisting}[language=c]
struct vifctl vc = { 0 };
vc.vifc_vifi = 1;
vc.vifc_flags = VIFF_USE_IFINDEX;
vc.vifc_lcl_ifindex = 2;

setsockopt(igmp_sock, IPPROTO_IP, MRT_ADD_VIF, &vc, sizeof(vc));
\end{lstlisting}

\noindent We are creating a VIF representing real network interface with index 2. This VIF will
get assigned the index 1. We can use an interface address to target a real
interface instead of interface index, or specify destination address for the
IPIP tunnel. If we set the \ttt{VIFF\_REGISTER} flag, special device will be created.

The kernel does not provide regular routing tables. The structure used for
forwarding is called Multicast Forwarding Cache (MFC). A MFC entry can be added
by filling a \ttt{struct mfcctl} and calling \ttt{MRT\_ADD\_MFC} or
\ttt{MRT\_ADD\_MFC\_PROXY}. The difference between these two will be
explained later.
\nomenclature{MFC}{Multicast Forwarding Cache}

\begin{lstlisting}[language=c]
struct mfcctl mc = { 0 };

mc.mfcc_origin = in_source;
mc.mfcc_mcastgrp = in_group;
mc.mfcc_parent = vifi;

for (int i = 0; i < MAXVIFS; i++)
    mc.mfcc_ttls[i] = ttl_threshold[i];

setsockopt(igmp_sock, IPPROTO_IP, MRT_ADD_MFC, &mc, sizeof(mc));
\end{lstlisting}

\noindent The basic meaning of an MFC entry is hidden in the field \ttt{mfcc\_ttls}. It is
an array of \ttt{MAXVIFS} integers, indexed by the VIF index. Whenever a field for the VIF is
either 0 or 255, we call the VIF inactive for this entry, otherwise we call it
active.

The entries are indexed by source and group address. When a packet for a~pair $(S,G)$
arrives on any interface, it is first checked whether it came from the
interface \ttt{mfcc\_parent}. If not, it is dropped immediately. Otherwise
for every active VIF~$i$, the packet's TTL is compared to the threshold in
\ttt{mfcc\_ttls[$i$]}. If the packet's TTL is greater, it is copied to the
VIF~$i$.

In Linux, there are two more types of MFC entries: $(*,G)$ and $(*,*)$. These
are added by a call \ttt{MFC\_ADD\_MFC\_PROXY}, and they have a completely different
meaning. I haven't found any explanation for these, my assumptions are based on
studying the kernel source code \cite{linux}. The proxy entries are used to build a shared
tree, possibly statically, common for all groups. First, you have to choose one
uplink interface~$i$, and add a~$(*,*)$ entry with a \ttt{mfcc\_parent} set
to~$i$. Whenever a packet comes to an active interface, and no $(*, G)$ MFC
entry is present yet, it is forwarded to the uplink.

For a simple example, imagine your home router. It probably has three
interfaces: the uplink \ttt{wan}, wired \ttt{lan} and wireless
\ttt{wlan} downlinks.\footnote{Those two downlinks are probably internally
bridged in your router, but not in this example.} A $(*,*)$ entry with
\ttt{wan} and \ttt{lan} as active interfaces would enable you to send
multicast packets from hosts on \ttt{lan}, which would then be forwarded to
\ttt{wan}. Multicast packets sent on \ttt{wlan} would be dropped.

Then, you can specify some~$(*,G)$ entries. Having a~$(*,G)$ entry with parent~$p$,
when a packet comes on an interface~$j$, and there is a~$(*,*)$ entry with
both~$p$ and~$j$ active, the packet is forwarded to all interfaces active in
the~$(*,G)$ entry with the exception of~$j$. Following the previous example,
you would add an~$(*,G)$ entry with all three interfaces active for a group
that you want to be delivered to both of your links. Any multicast packet
received on either \ttt{wan} or \ttt{lan} would be forwarded to the
remaining interfaces.

Note that even when a $(*,G)$ entry is present, the kernel looks for a $(*,*)$
entry to filter the incoming interface. There can be more $(*,*)$ entries with
different uplinks, but it will work only when the active interfaces are
disjoint. There are no checks to prevent you from adding an entry that would
never be found, so it is hard to guess which behavior is intentional and which
is not.

You may wonder whether it is possible to add unicast routes for multicast
addresses. It is, but their meaning is different. These routes are never used
to forward packets. Instead, when you open a socket and do not specify an
outgoing interface for a multicast packet (\ttt{IP\_MULTICAST\_IF}), the interface
will be selected according to these routes. If no matching route exists, the sending
will fail.

\chapter{Protocol Independent Multicast}

Unicast routing protocols try to cooperatively build a tree for every
destination network. This tree is a spanning tree over some graph of networks,
and is rooted at its destination. Every packet travels up the tree until
reaching its destination.

Multicast packets have no single destination, so multicast routing protocols
have to take slightly different approach. For every multicast group they still
try to build a tree connecting all sources and listeners of traffic. Members of
Protocol Independent Multicast (PIM) routing protocols use multicast routing
information base (MRIB), which is often derived from the unicast routing table.

In this area BIRD can take advantage, because it is designed to run multiple
routing protocols at a time and share information between them. One can for
example route unicast traffic using BGP and use OSPF with different metric to
supply completely different topology to multicast routing.

Every group managed by PIM has its \emph{Rendezvous-point} address (RPA)
associated. Assigning it is a matter of local configuration. RP is a router
specially configured for this purpose. Unicast delivery tree for RPA is then
called RP tree. With respect to this tree we can define a few more terms:

\emph{Rendezvous-point link} is simply a link where RP address belongs.
\emph{Upstream interface} for every router is the next-hop interface towards
RP, and \emph{upstream neighbor} is the next-hop router.

While unicast packet at first hops over a number of routers, then it is sent to
its destination, multicast packets are always sent on the link directly. This
way every router on link can hear it. In order to avoid duplicate forwards, on
every link other than RPL election takes place. Winning router is called
\emph{designated router}. This designated router is responsible for
distribution of information taken from IGMP messages to upstream neighbor. This
is made by PIM Join/Prune messages. This mechanism is almost reimplementation
of IGMPv2 in PIM context. When there is a host on link which joined group $G$
by IGMP, or any router that joined $G$ by PIM, DR joins $G$ on its upstream.

Let me first explain briefly how PIM Sparse mode works. Suppose there already
is a tree built for group $G$ with some listeners. Suddenly node $S$ starts
sending multicast packets onto link. Then DR on this link ($A$) takes these
packets and encapsulates them into ordinary unicast packets. These are then
sent to RP with special type of PIM messages. This is done actively by the
routing daemon. RP decapsulates them, and temporarily distributes them down the
tree. Routers are responsible for picking packets coming on upstream and
forwarding them down on links where they are acting as designated routers.

In the meantime, $A$ communicates with RP through some routers in between. They
all establish a multicast routed path exclusively for this source. Then
encapsulating is stopped and the traffic can flow as usual. This complicated
reactive behaviour is needed, because building a distribution tree in advance
for every group and every possible source would be too expensive.

Furthermore, RP has to be an existing router, and must be running and have
enough processing power to decapsulate traffic from all sources until
muliticast forwarding is established. Bidirectional PIM takes different
approach.

\section{Bidirectional PIM}

Everything is made a lot simpler, when we do not expect hosts joining specific
sources. Routers can have a lot smaller routing tables, and even build them
proactively. No encapsulating is then needed and processing power is saved. So
first clarify things in terms of PIM-BIDIR.

Designated router is called designated forwarder (DF) and takes additional
responsibilities. There are differences in how the election works. RP can be
just unicast routable address, and is used just as virtual root of the
distribution tree. Every multicast packet from any source than travels up the
tree, being forwarded by DFs, and then down to branches having listeners.

\subsection{Designated forwarder election}

Every router running PIM-BIDIR knows from MRIB what its upstream interface is,
and its metric towards RP. Router cannot be DF on its upstream link, so
whenewer it sends offer, advertises infinite metric. Election never takes place
on RPL, because it is the upstream link for every router.

There are four types of messages in the election: Offer, Winner, Backoff and
Pass. Routers advertise their metrics as pairs of (protocol preference, distance).
State machine with all details is described in \rfc{5051}. The main idea of the
election is that every router advertises it's metric in Offer messages. When
a router hears better Offer, it remains silent for a while. If not, after
advertising its metric three times, it claims itself a winner. Winning is
announced by Winner messages, and every other router notes sender as the
winner and remembers its metric.

Two noteworthy events may happen. If winner changes its metric to worse value,
it reannounces itself as a winner. When any other router has better metric,
replies with an Offer immediately. Similarly, when losing router changes its
metric to be better than winner's, it sends an Offer.

When metric changes at one router, it is often reaction to some router losing
link, and means this change will propagate to other routers, changing their
metric too. In order to spare bandwidth, current winner reacts to better offer
with a Backoff message. It instructs target router to wait a while until
unicast routing stabilizes. After a short period without other Offers, winner
passes DF role in Pass message.

\TODO{Pathological situations}

\subsection{Packet forwarding}

Suppose source $S$ starts transmitting packets on its link $L_0$. Then
currently acting DF will forward them immediately to its upstream link $L_1$.
DF on $L_1$ forwards it to its upstream, $L_2$, and so on. This is repeated
until the packet reaches RPL. As there is no DF on RPL, forwarding up stops
there.

Simultaneously, all routers listen for packets on their upstream links, and
forwards packets from there to links where they are acting DFs. This way packet
is forwarded from $L_0, L_1, \dots$ in reverse direction -- downstream.
Downstream forwarding is conditioned by having listeners active. Upstream
forwarding needs to be done even when there are no listeners.

You can see that every DF can forward packets for one group upstream and
downstream at the same time, explaining the name "Bidirectional".


\chapter*{Conclusion}
\addcontentsline{toc}{chapter}{Conclusion}

We have succeeded in implementing the multicast routing features in the BIRD.
Because the protocol is fairly complex, it deserves to be more tested and, what
is more important, reviewed by the BIRD team. They were too busy to deep-study
my implementation continuously. But merging the PIM into BIRD is on the plan.

The source code for the whole project is available as a Git repository,
located at the CZ.NIC GitLab:

\centerline{\ttt{https://gitlab.labs.nic.cz/Aearsis/bird.git}}

In the repository, you can find my work on the branch \ttt{pim}. Or, the
current snapshot is attached as a~patchset attached to the electronic version
of the thesis.

There was a considerable amount of obstacles during the work. First, I had to
learn how the BIRD itself works. It is not a small project, and there are lots
of internal details one must know before starting to contribute in a major way.
Also, I have decided to change the existing BIRD code as little as possible and
write my code in the spirit of the existing one, to follow the coding practices
and used design patterns. That made me study the code for even longer time, but
hopefully it will pay off when merging my code to the main development branch.

Next, there were problems with undocumented code in Linux kernel. Looking back,
when we know where to look for the issues, they seem trivial. But they blocked
the work for a while. Namely, all the surprises around the IGMP control socket,
details about how MFC entries work, or the Linux bridge doing IGMP snooping by
default.

Naturally, there are many ways how to expand this project. The first direction
should be support for IPv6. It should be fairly simple, because the internal
structures are prepared for it. Also, the PIM itself should be able to run on
IPv6 out of the box, but this behavior is not tested. One just needs to
implement the MLD protocol (IPv6 equivalent of the IGMP), and the IPv6 kernel
route synchronization. In the kernel, the v6 API looks like a copy of the v4,
on which a few text substitutions were applied.

Next, it would be great to implement another modes of PIM, especially the PIM-SM mode.
Again, the design of the PIM protocol is prepared for it. One would need to
add the different types of packets PIM-SM uses and prepend source address to
the internal structures. It will probably not be that easy, but we do not see
any fundamental problem now.

For the bravest, the part that needs reimplementation the most is the Linux
kernel itself. The implementation of multicast forwarding is copied from BSD to
stay compatible with the mrouted daemon, and it reflects how mrouted works.
The API should better follow the structure of the protocols instead of the
structure of a particular routing daemon. Also, the multicast input, IGMP
processing and other related chunks of code looks more like a hack than
a systematic solution. An obvious solution is to use netlink, through which the
rest of routing is already controlled.


%%% Bibliography
%%% Bibliography (literature used as a source)
%%%
%%% We employ bibTeX to construct the bibliography. It processes
%%% citations in the text (e.g., the \cite{...} macro) and looks up
%%% relevant entries in the bibliography.bib file.
%%%
%%% The \bibliographystyle command selects, which style will be used
%%% for references from the text. The argument in curly brackets is
%%% the name of the corresponding style file (*.bst). Both styles
%%% mentioned in this template are included in LaTeX distributions.

% \bibliographystyle{plainnat}    %% Author (year)
\bibliographystyle{unsrt}     %% [number]

\renewcommand{\bibname}{Bibliography}

%%% Generate the bibliography. Beware that if you cited no works,
%%% the empty list will be omitted completely.

\bibliography{rfc}

%%% If case you prefer to write the bibliography manually (without bibTeX),
%%% you can use the following. Please follow the ISO 690 standard and
%%% citation conventions of your field of research.

\begin{thebibliography}{99}


\end{thebibliography}


%%% Figures used in the thesis (consider if this is needed)
%%% \listoffigures

%%% Tables used in the thesis (consider if this is needed)
%%% In mathematical theses, it could be better to move the list of tables to the beginning of the thesis.
%%% \listoftables

%%% Abbreviations used in the thesis, if any, including their explanation
%%% In mathematical theses, it could be better to move the list of abbreviations to the beginning of the thesis.
\chapwithtoc{List of Abbreviations}

%%% Attachments to the bachelor thesis, if any. Each attachment must be
%%% referred to at least once from the text of the thesis. Attachments
%%% are numbered.
%%%
%%% The printed version should preferably contain attachments, which can be
%%% read (additional tables and charts, supplementary text, examples of
%%% program output, etc.). The electronic version is more suited for attachments
%%% which will likely be used in an electronic form rather than read (program
%%% source code, data files, interactive charts, etc.). Electronic attachments
%%% should be uploaded to SIS and optionally also included in the thesis on a~CD/DVD.
\chapwithtoc{Attachments}

\openright
\end{document}
